% Wersja źródłowa: audyt III, jednostka 02
\bigskip
\setcounter{section}{1}
\section{Podprzestrzenie, produkty, ilorazy i topologie uniwersalne}\label{sec:konstrukcje}

\subsection{Dlaczego potrzebujemy konstrukcji nowych przestrzeni?}

Z jednej przestrzeni liniowo-topologicznej (TVS) powstają naturalnie nowe:
zbiór funkcji spełniających liniowy warunek, rodzina funkcji zależnych
od wielu parametrów, przestrzeń funkcji rozpatrywanych z dokładnością
do pewnej relacji równoważności. Odpowiadają im odpowiednio
\emph{podprzestrzenie}, \emph{produkty} oraz \emph{ilorazy}.
Istotą konstrukcji jest nie tylko zdefiniowanie zbioru i działań,
lecz także dobranie topologii tak, aby właściwe odwzorowania były
ciągłe. W dalszych rozdziałach ta idea posłuży do budowy
przestrzeni funkcji testowych i dystrybucji.

Przyjmujemy konwencje \hyperref[sec:tvs]{rozdziału 1}: $\K\in\{\R,\C\}$,
a wszystkie przestrzenie liniowo-topologiczne są Hausdorffa,
chyba że wyraźnie zaznaczymy inaczej. Pojęcie \emph{topologii
indukowanej} i \emph{topologii produktowej} rozumiemy w zwykłym sensie
topologii ogólnej. W tej części przypomnimy ich precyzyjne własności.

\subsection{Podprzestrzenie liniowo-topologiczne}

\begin{twierdzenie}[Topologia podprzestrzeni]\label{thm:subspace-tvs}
Niech $X$ będzie przestrzenią liniowo-topologiczną nad $\K$,
a $M\subseteq X$ podprzestrzenią liniową. Z topologią indukowaną
$\tau_M=\{O\cap M:O\subseteq X\text{ jest otwarty}\}$
przestrzeń $M$ jest przestrzenią liniowo-topologiczną nad $\K$.
Jeżeli $X$ jest Hausdorffa, to również $M$ jest Hausdorffa.
Otoczenia zera w $M$ mają bazę postaci $U\cap M$, gdzie
$U$ przebiega otoczenia zera w $X$.
\end{twierdzenie}
\begin{proof}
Ograniczenie dodawania $+_X$ do $M\times M$ przyjmuje wartości
w $M$, a topologia na $M\times M$ jest topologią podprzestrzeni
$X\times X$. Jeżeli $O\cap M$ jest otwarty w $M$, to
\[
 (+_M)^{-1}(O\cap M)=(+_X)^{-1}(O)\cap(M\times M),
\]
więc dodawanie w $M$ jest ciągłe. Identyczny argument dla
$\K\times M\subseteq\K\times X$ daje ciągłość mnożenia
przez skalar. Dziedziczenie własności Hausdorffa jest własnością
każdej podprzestrzeni topologicznej. Ostatnie zdanie wynika
wprost z definicji topologii indukowanej i faktu, że $0\in M$.
\end{proof}

\begin{lemat}[Domknięcie podprzestrzeni]\label{lem:closure-subspace}
Niech $M$ będzie podprzestrzenią liniową przestrzeni
liniowo-topologicznej $X$, a $\mathcal N_X(0)$ rodziną
otoczeń zera w $X$. Wówczas $\overline M$ jest podprzestrzenią
liniową oraz
\[
 \overline M=\bigcap_{U\in\mathcal N_X(0)}(M+U),
 \qquad M+U=\{m+u:m\in M,\ u\in U\}.
\]
\end{lemat}
\begin{proof}
Jeżeli $x,y\in\overline M$, to dla dowolnego otoczenia $U$
zera wybieramy symetryczne otoczenie $V$ z $V+V\subseteq U$.
Istnieją $m_x,m_y\in M$ takie, że $x-m_x\in V$ oraz
$y-m_y\in V$. Stąd $(x+y)-(m_x+m_y)\in U$,
więc $x+y\in\overline M$. Dla $\lambda\ne0$ ciągłość
mnożenia przez skalar (a nawet fakt, że jest homeomorfizmem)
daje $\lambda\overline M=\overline{\lambda M}=\overline M$;
przypadek $\lambda=0$ jest oczywisty.

Jeżeli $x\in\overline M$, to dla każdego $U$ możemy wybrać
symetryczne $V\subseteq U$; przecięcie $(x+V)\cap M$
jest niepuste, więc $x=m-v\in M+U$ dla odpowiednich
$m\in M$, $v\in V$. Odwrotnie, niech $x\in\bigcap_U(M+U)$.
Dla otoczenia $x+V$ punktu $x$ wybieramy symetryczne $V$;
z założenia $x=m+v$ z $m\in M$, $v\in V$,
a więc $m=x-v\in x+V$. Każde otoczenie $x$ spotyka $M$,
czyli $x\in\overline M$.
\end{proof}

\begin{przyklad}[Dwie topologie na wielomianach]
Niech $P=\K[t]$, a $X=C([0,1],\K)$ z normą supremum.
Topologia indukowana na $P$ jest unormowana, lecz $P$
nie jest domknięte w $X$: wielomiany są gęste w $X$
(z twierdzenia Weierstrassa o aproksymacji). W szczególności
zupełność $X$ nie przechodzi na dowolną podprzestrzeń.
W przyszłości udowodnimy, że podprzestrzeń przestrzeni Banacha
jest zupełna w normie indukowanej wtedy i tylko wtedy, gdy
jest domknięta.
\end{przyklad}

\begin{stwierdzenie}[Uniwersalna własność podprzestrzeni]
Niech $M\subseteq X$ będzie podprzestrzenią liniową przestrzeni
liniowo-topologicznej $X$, $i:M\hookrightarrow X$ naturalnym
włożeniem, $Z$ dowolną przestrzenią topologiczną oraz
$f:Z\to M$ odwzorowaniem. Wówczas $f$ jest ciągłe
wtedy i tylko wtedy, gdy $i\circ f:Z\to X$ jest ciągłe.
\end{stwierdzenie}
\begin{proof}
Konieczność jest natychmiastowa, gdyż $i$ jest ciągłe.
Jeśli $i\circ f$ jest ciągłe, to dla otwartego $O\cap M$
w $M$ mamy $f^{-1}(O\cap M)=(i\circ f)^{-1}(O)$,
więc $f$ jest ciągłe.
\end{proof}

\subsection{Produkty dowolnych rodzin przestrzeni}

\begin{definicja}[Produkt przestrzeni liniowo-topologicznych]
Niech $I$ będzie zbiorem indeksów, a $(X_i)_{i\in I}$ rodziną
przestrzeni liniowo-topologicznych nad $\K$. Na przestrzeni
liniowej $X=\prod_{i\in I}X_i$ działań współrzędnych
wprowadzamy \emph{topologię produktową}: jej bazę tworzą zbiory
$\prod_{i\in I}O_i$, gdzie każde $O_i\subseteq X_i$ jest
otwarte i $O_i=X_i$ dla wszystkich poza skończenie wieloma
indeksami. Rzut na $i$-tą współrzędną oznaczamy przez
$\pi_i:X\to X_i$.
\end{definicja}

\begin{twierdzenie}[Produkt i jego własność uniwersalna]\label{thm:product-tvs}
Niech $(X_i)_{i\in I}$ będzie rodziną przestrzeni
liniowo-topologicznych Hausdorffa nad $\K$.
Wówczas $X=\prod_{i\in I}X_i$ z topologią produktową
jest przestrzenią liniowo-topologiczną Hausdorffa, a każdy
$\pi_i$ jest ciągłym, liniowym i otwartym odwzorowaniem na $X_i$.
Dla dowolnej przestrzeni topologicznej $Z$ odwzorowanie
$f:Z\to X$ jest ciągłe wtedy i tylko wtedy, gdy
$\pi_i\circ f:Z\to X_i$ jest ciągłe dla każdego $i\in I$.
\end{twierdzenie}
\begin{proof}
Rzuty są ciągłe z definicji topologii produktowej. Dla zbioru
bazowego $\prod_i O_i$ zachodzi $\pi_j(\prod_i O_i)=O_j$,
ponieważ wszystkie pozostałe $O_i$ są niepuste, o ile
rozpatrujemy niepusty zbiór bazowy; rzut zbioru pustego
jest pusty. Zatem rzuty są otwarte.

Dodawanie $a:X\times X\to X$ jest ciągłe, bo każde
$\pi_i\circ a:(x,y)\mapsto x_i+y_i$ jest ciągłe
(jest to złożenie rzutowań i dodawania w $X_i$).
To samo rozumowanie stosuje się do mnożenia skalarnego
$m:\K\times X\to X$. Przestrzeń $X$ jest Hausdorffa,
ponieważ dwa różne elementy różnią się na pewnej współrzędnej,
a tę można rozdzielić otoczeniami w odpowiednim $X_i$.

Jeśli wszystkie $f_i=\pi_i\circ f$ są ciągłe,
to dla podzbioru bazowego $\prod_iO_i$ mamy
$f^{-1}(\prod_iO_i)=\bigcap_{i\in F}f_i^{-1}(O_i)$,
gdzie $F=\{i:O_i\ne X_i\}$ jest skończony;
jest to zbiór otwarty. Implikacja odwrotna wynika
ze złożenia odwzorowań ciągłych.
\end{proof}

\begin{uwaga}[Dlaczego tylko skończenie wiele warunków?]
W produkcie nieskończonym wymaganie, aby wszystkie
współrzędne jednocześnie należały do wybranych małych otoczeń,
może zniszczyć ciągłość działań. Na przykład w $\K^{\N}$
zbiory $\{x:|x_n|<1\text{ dla każdego }n\}$ nie są otwarte
w topologii produktowej. Ograniczanie skończenie wielu
współrzędnych jest kluczowe dla uniwersalnej własności produktu.
\end{uwaga}

\begin{twierdzenie}[Ograniczoność w produkcie]\label{thm:product-bounded}
Niech $(X_i)_{i\in I}$ będzie rodziną przestrzeni
liniowo-topologicznych i niech $B\subseteq\prod_{i\in I}X_i$.
Wówczas $B$ jest ograniczony w sensie von Neumanna wtedy
i tylko wtedy, gdy $\pi_i(B)$ jest ograniczony w $X_i$
dla każdego $i\in I$.
\end{twierdzenie}
\begin{proof}
W jedną stronę stosujemy rezultat z \hyperref[sec:tvs]{rozdziału 1}: ciągły operator liniowy
przenosi zbiory ograniczone na zbiory ograniczone. Stosujemy go do ciągłych rzutów $\pi_i$.
Odwrotnie, niech $U$ będzie otoczeniem zera w produkcie.
Możemy wybrać bazowe otoczenie $V=\prod_i V_i\subseteq U$,
gdzie $V_i=X_i$ poza skończonym zbiorem $F$, a dla $i\in F$
przyjmijmy $V_i$ zrównoważone. Dla każdego $i\in F$ istnieje
$t_i>0$ takie, że $\pi_i(B)\subseteq t_iV_i$.
Jeśli $F$ jest niepusty, ustalmy $t\geq\max(1,\max_{i\in F}t_i)$.
Zrównoważenie $V_i$ daje $t_iV_i\subseteq tV_i$, więc
$B\subseteq tV\subseteq tU$. Jeżeli $F$ jest pusty,
wystarczy przyjąć $t=1$.
\end{proof}

\begin{przyklad}[Dwa rodzaje zbieżności funkcji]
Niech $X=C([0,1],\K)$. Norma supremum wyznacza topologię
zbieżności jednostajnej; topologię zbieżności punktowej
otrzymujemy przez zanurzenie $X\subseteq\K^{[0,1]}$.
Identyczność z topologii normowej do punktowej jest ciągła.
Nie jest jednak homeomorfizmem: funkcje
$f_n(t)=\max\{0,1-n^2|t-1/n|\}$, $n\geq2$, są ciągłe,
zbiegają punktowo do zera, ale $\|f_n\|_\infty=1$.
Istotnie, dla każdego ustalonego $t\in[0,1]$ nośnik $f_n$
nie zawiera $t$ dla wszystkich dostatecznie dużych $n$.
\end{przyklad}

\begin{twierdzenie}[Funkcjonały na produkcie przestrzeni skalarnych]\label{thm:product-dual}
Niech $I$ będzie dowolnym zbiorem, a $X=\K^I$ wyposażmy
w topologię produktową. Każdy ciągły funkcjonał liniowy
$L:X\to\K$ ma postać
\[
 L(x)=\sum_{i\in F}a_i x_i
\]
przy pewnym skończonym zbiorze $F\subseteq I$ oraz
skalarach $a_i\in\K$. Odwrotnie, każdy taki funkcjonał
jest ciągły.
\end{twierdzenie}
\begin{proof}
Z ciągłości $L$ w zerze istnieje bazowe otoczenie zera
$U=\{x:|x_i|<\varepsilon_i\text{ dla }i\in F\}$
takie, że $|L(x)|<1$ dla $x\in U$.
Jeśli $x_i=0$ dla każdego $i\in F$, to $\lambda x\in U$
dla dowolnego $\lambda\in\K$. Stąd
$|\lambda L(x)|<1$ dla wszystkich $\lambda$, więc $L(x)=0$.
Dla dowolnego $x\in X$ odejmijmy skończoną sumę
$\sum_{i\in F}x_i e_i$, gdzie $e_i$ oznacza funkcję
równą $1$ na $i$ i $0$ poza nim. Pozostały wektor
ma zera na $F$, dlatego $L(x)=\sum_{i\in F}L(e_i)x_i$.
Odwrotność wynika z ciągłości rzutów i działań skalarnych.
\end{proof}

\begin{uwaga}
Twierdzenie~\ref{thm:product-dual} jest pierwszą wyraźną
manifestacją tego, że \emph{przestrzeń ciągłych funkcjonałów liniowych}
może być znacznie mniejsza od przestrzeni funkcjonałów algebraicznych. Nie należy go
mylić z analogiczną własnością dowolnych produktów:
ciągły operator $\K^{\N}\to\K^{\N}$ może zależeć
od nieskończenie wielu współrzędnych (np. identyczność).
\end{uwaga}

\subsection{Przestrzeń ilorazowa i topologia ilorazowa}

\begin{definicja}[Iloraz przez podprzestrzeń]
Niech $X$ będzie przestrzenią liniowo-topologiczną nad $\K$,
a $M\subseteq X$ podprzestrzenią liniową. Na przestrzeni
wektorowej $X/M=\{x+M:x\in X\}$ wprowadzamy
\emph{topologię ilorazową}: zbiór $W\subseteq X/M$
jest otwarty wtedy i tylko wtedy, gdy $q^{-1}(W)$
jest otwarty w $X$, gdzie $q:X\to X/M$, $q(x)=x+M$.
\end{definicja}

\begin{twierdzenie}[Iloraz jako TVS]\label{thm:quotient-tvs}
Niech $X$ będzie dowolną przestrzenią liniowo-topologiczną,
niekoniecznie Hausdorffa, a $M\subseteq X$ podprzestrzenią.
Wówczas:
\begin{enumerate}[label=\textup{(\roman*)}]
\item odwzorowanie $q:X\to X/M$ jest liniowe, ciągłe,
      surjektywne i otwarte;
\item topologia ilorazowa czyni $X/M$ przestrzenią
      liniowo-topologiczną;
\item $X/M$ jest Hausdorffa wtedy i tylko wtedy,
      gdy $M$ jest domknięta w $X$.
\end{enumerate}
\end{twierdzenie}
\begin{proof}
Liniowość i surjektywność $q$ wynikają z konstrukcji
ilorazu, a ciągłość wynika z definicji topologii ilorazowej.
Jeżeli $O\subseteq X$ jest otwarty, to
\[
 q^{-1}(q(O))=O+M=\bigcup_{m\in M}(O+m)
\]
jest otwarty, ponieważ każde przesunięcie $O+m$
jest otwarte. Zatem $q(O)$ jest otwarty w $X/M$.

Wykażemy ciągłość dodawania w ilorazie. Odwzorowanie
$q\times q:X\times X\to(X/M)\times(X/M)$ jest ciągłą,
otwartą surjekcją: otwartość wynika z faktu, że obrazy
zbiorów bazowych $O_1\times O_2$ są produktami
otwartych $q(O_1)\times q(O_2)$. Każda otwarta surjekcja
ciągła jest odwzorowaniem ilorazowym. Mamy równość
$+_{X/M}\circ(q\times q)=q\circ+_X$, a prawa strona
jest ciągła. Z własności odwzorowania ilorazowego wynika
ciągłość $+_{X/M}$. Podobnie odwzorowanie
$\id_\K\times q:\K\times X\to\K\times(X/M)$ jest
otwartą surjekcją, więc z równości
$m_{X/M}\circ(\id_\K\times q)=q\circ m_X$
wynika ciągłość mnożenia skalarnego. Tym samym $X/M$
jest TVS, także gdy $M$ nie jest domknięta.

W dowolnej TVS singleton $\{0\}$ jest domknięty
wtedy i tylko wtedy, gdy przestrzeń jest Hausdorffa
(por. twierdzenie o domknięciu zera z \hyperref[sec:tvs]{rozdziału 1}).
Ponieważ $q^{-1}(\{0+M\})=M$, definicja topologii
ilorazowej pokazuje, że $\{0+M\}$ jest domknięty
w $X/M$ wtedy i tylko wtedy, gdy $M$ jest domknięta w $X$.
\end{proof}

\begin{uwaga}[Otoczenia w ilorazie]\label{rem:q-neighborhood}
Ponieważ $q$ jest otwarty, zbiory $q(U)$, gdzie
$U\in\mathcal N_X(0)$, tworzą bazę otoczeń zera
w $X/M$. Dokładniej: dla każdego otoczenia $W$
punktu $0+M$ zbiór $q^{-1}(W)$ jest otoczeniem zera,
a $q(q^{-1}(W))=W$; z drugiej strony każdy $q(U)$
jest otoczeniem zera, bo zawiera $q(O)$ dla pewnego
otwartego otoczenia $O\subseteq U$ zera.
\end{uwaga}

\begin{center}
\begin{tikzpicture}[x=0.95cm,y=0.82cm,>=Latex]
\draw[->] (-0.5,0)--(5.35,0) node[right] {$x$};
\draw[->] (0,-0.55)--(0,3.75) node[above] {$y$};
\foreach \yy in {0.8,1.7,2.65} {
  \draw[blue!55!black,thick] (0.2,\yy)--(4.7,\yy);
}
\filldraw[blue!55!black] (1.0,1.7) circle (1.7pt);
\filldraw[blue!55!black] (3.3,1.7) circle (1.7pt);
\node[anchor=west] at (4.75,1.7) {$x_0+M$};
\node[below] at (2.4,-0.10) {$M=\R\times\{0\}$};
\draw[->,very thick] (5.15,2.8) -- (6.4,2.8)
  node[midway,above] {$q$};
\draw[->] (7.2,-0.5)--(7.2,3.65);
\foreach \yy in {0.8,1.7,2.65} {
 \filldraw[blue!55!black] (7.2,\yy) circle (2.1pt);
}
\node[right] at (7.4,1.7) {$[x_0]$};
\node[below] at (7.2,-0.65) {$\R^2/M\simeq\R$};
\end{tikzpicture}

{\small Rysunek 2.1. Iloraz $\R^2/M$, gdzie $M$ jest osią poziomą.
Każda pozioma prosta jest jedną klasą równoważności
(prawa strona przedstawia zbiór klas).}
\end{center}

\begin{twierdzenie}[Własność uniwersalna ilorazu]\label{thm:quotient-universal}
Niech $X$ będzie TVS nad $\K$, $M\subseteq X$ podprzestrzenią,
a $q:X\to X/M$ rzutowaniem. Niech $Z$ będzie dowolną
przestrzenią topologiczną, a $g:X/M\to Z$ odwzorowaniem.
Wówczas $g$ jest ciągłe wtedy i tylko wtedy,
gdy $g\circ q$ jest ciągłe. Ponadto jeśli $Z$ jest TVS,
a $T:X\to Z$ jest liniowe i $M\subseteq\ker T$, to
istnieje dokładnie jedno liniowe $\widetilde T:X/M\to Z$
takie, że $T=\widetilde T\circ q$; jest ono ciągłe
dokładnie wtedy, gdy $T$ jest ciągłe.
\end{twierdzenie}
\begin{proof}
Pierwsza równoważność wynika z definicji topologii ilorazowej:
$g^{-1}(O)$ jest otwarty w $X/M$ dokładnie wtedy,
gdy $(g\circ q)^{-1}(O)$ jest otwarty w $X$, dla każdego
otwartego $O\subseteq Z$. Jeśli $M\subseteq\ker T$,
to wzór $\widetilde T(x+M)=T(x)$ jest dobrze określony:
z równości $x+M=y+M$ wynika $x-y\in M$, czyli $T(x)=T(y)$.
Wzór jest liniowy, a surjektywność $q$ zapewnia jednoznaczność.
Równoważność ciągłości stosujemy do $g=\widetilde T$.
\end{proof}

\begin{wniosek}[Pierwsze twierdzenie o izomorfizmie --- uwaga topologiczna]
Niech $T:X\to Y$ będzie ciągłym, liniowym odwzorowaniem
między przestrzeniami liniowo-topologicznymi Hausdorffa.
Wówczas odwzorowanie
$\widetilde T:X/\ker T\to T(X)$, $x+\ker T\mapsto T(x)$,
jest ciągłą bijekcją liniową, gdzie $T(X)$ ma topologię
indukowaną z $Y$. Jest izomorfizmem liniowo-topologicznym
wtedy i tylko wtedy, gdy $T:X\to T(X)$ jest odwzorowaniem
otwartym (przy topologii podprzestrzeni na obrazie).
\end{wniosek}
\begin{proof}
Jądro $T$ jest domknięte, ponieważ $T$ jest ciągłe,
a $Y$ jest Hausdorffa; dlatego iloraz jest Hausdorffa.
Istnienie i ciągłość $\widetilde T$ daje
\ref{thm:quotient-universal}; jest ona różnowartościowa,
bo $\ker T$ jest dokładnie jądrem $T$.
Jeśli $T$ jest otwarte na swój obraz, to
$\widetilde T(q(O))=T(O)$ jest otwarty dla każdego otwartego
$O\subseteq X$, a z otwartości $q$ i opisu zbiorów
otwartych w ilorazie wynika otwartość $\widetilde T$.
Odwrotnie, jeżeli $\widetilde T$ jest homeomorfizmem,
to $T=\widetilde T\circ q$ jest otwarte na obraz jako
złożenie dwóch odwzorowań otwartych.
\end{proof}

\begin{przyklad}[Dlaczego ciągła bijekcja nie wystarcza]
Niech $X=C([0,1],\R)$ z normą supremum, a $Y$ to ta sama
przestrzeń liniowa z normą $\|f\|_1=\int_{[0,1]}|f(t)|\,\mathrm d\lambda(t)$.
Identyczność $T:X\to Y$ jest ciągłą bijekcją liniową,
ale nie jest otwarta. Jej jądro jest zerowe, a mimo to
$X/\ker T\simeq X$ nie jest topologicznie izomorficzne z $Y$
poprzez odwzorowanie indukowane przez $T$. Dowód braku
ciągłości odwrotności podaliśmy w \hyperref[sec:tvs]{rozdziale 1}.
\end{przyklad}

\begin{przyklad}[Iloraz przez jądro wartościowania]
Niech $X=C([0,1],\K)$ z normą supremum, a
$E_0(f)=f(0)$. Jądro $M=\{f:f(0)=0\}$ jest domknięte.
Rzut $X/M\to\K$, $[f]\mapsto f(0)$, jest izomorfizmem
liniowo-topologicznym. Istotnie, dla $a\in\K$ stała
funkcja $s(a)(t)=a$ jest ciągłym liniowym prawym odwrotnym
do $E_0$, a odwrotne odwzorowanie na iloraz ma postać
$a\mapsto q(s(a))$. Jest więc ciągłe jako złożenie ciągłych.
\end{przyklad}

\subsection{Norma ilorazowa: ważny przypadek szczególny}

\begin{twierdzenie}[Norma w ilorazie]\label{thm:quotient-norm}
Niech $(X,\|\cdot\|)$ będzie przestrzenią unormowaną,
a $M\subseteq X$ domkniętą podprzestrzenią liniową. Wzór
\[
 \|x+M\|_{X/M}:=\inf_{m\in M}\|x-m\|=\operatorname{dist}(x,M)
\]
definiuje normę na $X/M$ generującą topologię ilorazową.
Dla każdego $r>0$ zachodzi
$q(B_X(0,r))=B_{X/M}(0,r)$.
Jeżeli dodatkowo $X$ jest przestrzenią Banacha,
to $X/M$ jest przestrzenią Banacha.
\end{twierdzenie}
\begin{proof}
Zastąpienie $x$ przez $x+m_0$, $m_0\in M$, nie zmienia
zbioru wartości $\|x+m_0-m\|$ dla $m\in M$, więc wzór
jest dobrze określony. Jednorodność wynika z liniowości $M$;
dla $\lambda=0$ obie strony są równe zeru. Dla dowolnych
$m_1,m_2\in M$ nierówność trójkąta w $X$ daje
$\|(x+y)-(m_1+m_2)\|\leq\|x-m_1\|+\|y-m_2\|$.
Biorąc infima, otrzymujemy nierówność trójkąta w ilorazie.
Jeśli $\operatorname{dist}(x,M)=0$, to $x\in\overline M=M$,
więc $x+M=0+M$. Zatem mamy normę.

Równość kul zachodzi, bo $\|x+M\|_{X/M}<r$ dokładnie
wtedy, gdy istnieje $m\in M$ z $\|x-m\|<r$.
Wobec otwartości $q$ wynika z tego, że obrazy kul
są otwarte. Dla dowolnego zbioru $W\subseteq X/M$
otwartego w topologii normowej jego przeciwobraz
$q^{-1}(W)$ jest otwarty, gdyż
$\|q(x)-q(y)\|_{X/M}\leq\|x-y\|$.
W drugą stronę, jeśli $W$ jest otwarty w topologii ilorazowej,
to dla każdego $[x]\in W$ istnieje $r>0$ takie, że
$B_X(x,r)\subseteq q^{-1}(W)$; wtedy
$B_{X/M}([x],r)=q(B_X(x,r))\subseteq W$.
Topologie są więc identyczne.

Niech teraz $X$ będzie zupełne, a $(u_n)$ ciągiem Cauchy'ego
w $X/M$. Wybierzmy podciąg $(u_{n_k})$ taki, że
$\|u_{n_{k+1}}-u_{n_k}\|_{X/M}<2^{-k-1}$.
Ustalmy przedstawiciela $x_1$ klasy $u_{n_1}$.
Indukcyjnie wybieramy przedstawiciela $x_{k+1}$ klasy
$u_{n_{k+1}}$ z $\|x_{k+1}-x_k\|<2^{-k}$.
Jest to możliwe, gdyż definicja normy ilorazowej
pozwala, po ustaleniu $x_k$, dobrać przedstawiciela
kolejnej klasy tak, aby odległość była dowolnie bliska
normie różnicy klas. Szereg odległości jest zbieżny,
więc $(x_k)$ jest ciągiem Cauchy'ego w $X$ i zbiega
do pewnego $x\in X$. Z ciągłości $q$ mamy
$u_{n_k}=q(x_k)\to q(x)$. Ponieważ cały $(u_n)$
jest ciągiem Cauchy'ego, również zbiega do $q(x)$.
\end{proof}

\begin{uwaga}
Gdy $M$ nie jest domknięta, ten sam wzór definiuje
jedynie seminormę: $\|x+M\|_{X/M}=0$ może zachodzić
dla niezerowej klasy. W szczególności taki wzór nie
może generować topologii Hausdorffa na $X/M$.
\end{uwaga}

\subsection{Topologie początkowe: najoszczędniejszy wybór topologii}

\begin{definicja}[Topologia początkowa]
Niech $S$ będzie zbiorem, a $f_i:S\to Y_i$, $i\in I$,
rodziną odwzorowań do przestrzeni topologicznych $Y_i$.
\emph{Topologią początkową} na $S$ względem $(f_i)$
nazywamy najmniejszą topologię, względem której wszystkie
$f_i$ są ciągłe. Jest ona generowana przez podbazę
$\{f_i^{-1}(O):i\in I,\ O\subseteq Y_i\text{ otwarty}\}$.
\end{definicja}

\begin{twierdzenie}[Początkowa topologia liniowa]\label{thm:initial-tvs}
Niech $X$ będzie przestrzenią liniową nad $\K$, $(Y_i)_{i\in I}$
rodziną przestrzeni liniowo-topologicznych nad $\K$,
a $T_i:X\to Y_i$ odwzorowaniami liniowymi.
Topologia początkowa względem $(T_i)$ czyni $X$
przestrzenią liniowo-topologiczną, niekoniecznie Hausdorffa.
Dla dowolnej przestrzeni topologicznej $Z$ odwzorowanie
$f:Z\to X$ jest ciągłe wtedy i tylko wtedy, gdy
$T_i\circ f$ jest ciągłe dla każdego $i\in I$.
Jeśli wszystkie $Y_i$ są Hausdorffa, to $X$ jest Hausdorffa
dokładnie wtedy, gdy $\bigcap_{i\in I}\ker T_i=\{0\}$.
\end{twierdzenie}
\begin{proof}
Ciągłość odwzorowań $T_i$ wynika z definicji topologii.
Ponieważ $T_i\circ+_X=+_{Y_i}\circ(T_i\times T_i)$,
a odwzorowania $T_i\times T_i$ są ciągłe względem
topologii początkowej na $X$, wszystkie złożenia
$T_i\circ+_X$ są ciągłe. Z uniwersalnej własności
topologii początkowej wynika ciągłość dodawania na $X$.
Tak samo $T_i\circ m_X=m_{Y_i}\circ(\id_\K\times T_i)$
daje ciągłość mnożenia skalarnego. Równoważność dotycząca
$f$ wynika wprost z preobrazów podbazowych zbiorów otwartych.

Jeśli $x\ne y$ i istnieje $i$ z $T_i(x)\ne T_i(y)$,
to dzięki własności Hausdorffa $Y_i$ istnieją rozłączne
otoczenia tych obrazów, których przeciwobrazy rozdzielają
$x$ i $y$. Jeśli natomiast $0\ne z\in\bigcap_i\ker T_i$,
to każde bazowe otoczenie zera zawiera $z$, a więc
$X$ nie jest Hausdorffa na mocy kryterium Hausdorffa
z \hyperref[sec:tvs]{rozdziału 1}.
\end{proof}

\begin{przyklad}[Produkty i podprzestrzenie jako przypadki topologii początkowej]
Topologia produktowa na $\prod_iX_i$ jest topologią
początkową względem rzutów $\pi_i$.
Topologia podprzestrzeni na $M\subseteq X$ jest
początkową względem włożenia $M\hookrightarrow X$.
W szczególności naturalna topologia na $C([0,1])$
zbieżności punktowej pochodzi od rodziny funkcjonałów
wartościowania $f\mapsto f(t)$, $t\in[0,1]$.
\end{przyklad}

\subsection{Topologie końcowe: przypadek delikatniejszy}

\begin{definicja}[Topologia końcowa zbiorów]
Niech $S$ będzie zbiorem, a $f_i:Y_i\to S$, $i\in I$,
rodziną odwzorowań z przestrzeni topologicznych.
\emph{Topologią końcową} na $S$ nazywamy najsilniejszą
topologię, przy której wszystkie $f_i$ są ciągłe:
\[
 \tau_{\mathrm{fin}}=\{O\subseteq S:
    f_i^{-1}(O)\text{ jest otwarty w }Y_i
    \text{ dla wszystkich }i\in I\}.
\]
Dla przestrzeni ilorazowej $X/M$ otrzymujemy właśnie
tę topologię przy jednym odwzorowaniu $q:X\to X/M$.
\end{definicja}

\begin{uwaga}[Niebezpieczeństwo: topologia końcowa nie musi być liniowa]
Jeżeli $S$ jest przestrzenią liniową, a $f_i$ są liniowe,
topologia końcowa w kategorii \emph{przestrzeni topologicznych}
nie musi czynić $S$ przestrzenią liniowo-topologiczną.
Przykład poniżej pokazuje, że nie można bez dowodu
przenosić własności ilorazu na dowolne rodziny odwzorowań.
\end{uwaga}

\begin{przyklad}[Dwie osie i błędna topologia końcowa]\label{ex:final-cross}
Niech $S=\R^2$, $Y_1=Y_2=\R$ z topologiami zwykłymi,
a $f_1(t)=(t,0)$ i $f_2(t)=(0,t)$. W topologii końcowej
zbiór jednoelementowy $\{(1,1)\}$ jest otwarty,
gdyż oba jego przeciwobrazy są puste. Jednak $\{(0,0)\}$
nie jest otwarty: jego przeciwobraz przez $f_1$ jest
zbiorem $\{0\}\subseteq\R$. Przesunięcie o $(-1,-1)$
przenosi pierwszy zbiór na drugi, więc nie jest
homeomorfizmem. Tymczasem w każdej TVS wszystkie
przesunięcia są homeomorfizmami. A zatem ta topologia
końcowa nie jest topologią przestrzeni liniowo-topologicznej.
\end{przyklad}

\begin{definicja}[Końcowa topologia liniowo-topologiczna]
Niech $X$ będzie przestrzenią liniową nad $\K$, a
$T_i:X_i\to X$, $i\in I$, rodziną odwzorowań
liniowych z przestrzeni liniowo-topologicznych.
\emph{Końcową topologią liniowo-topologiczną}
nazywamy najsilniejszą topologię na $X$, która czyni $X$
przestrzenią liniowo-topologiczną i wszystkie $T_i$
odwzorowaniami ciągłymi. W tej definicji nie zakładamy
warunku Hausdorffa.
\end{definicja}

\begin{twierdzenie}[Istnienie końcowej topologii liniowo-topologicznej]
\label{thm:final-tvs}
Dla każdej przestrzeni liniowej $X$ nad $\K$ i rodziny
liniowych odwzorowań $T_i:X_i\to X$ z TVS istnieje
końcowa topologia liniowo-topologiczna. Ma ona następującą
własność: jeżeli $Z$ jest TVS, a $S:X\to Z$ odwzorowaniem
liniowym, to $S$ jest ciągłe dokładnie wtedy, gdy
$S\circ T_i$ jest ciągłe dla wszystkich $i\in I$.
\end{twierdzenie}
\begin{proof}
Rozpatrzmy rodzinę $\mathfrak T$ wszystkich topologii
liniowo-topologicznych na $X$, względem których wszystkie
$T_i$ są ciągłe. Nie jest ona pusta: zawiera topologię
trywialną. Niech $\tau$ będzie topologią generowaną
przez sumę wszystkich topologii z $\mathfrak T$.
Aby nie pomylić topologii na dziedzinach działań,
ustalmy otwarty w pewnej $\sigma\in\mathfrak T$ zbiór
$O\subseteq X$. Wówczas
$+_X^{-1}(O)$ jest otwarty w $(X,\sigma)\times(X,\sigma)$,
a więc również w $(X,\tau)\times(X,\tau)$.
Analogicznie $m_X^{-1}(O)$ jest otwarty w
$\K\times(X,\tau)$. Ponieważ takie zbiory $O$
tworzą podbazę $\tau$, oba działania są łącznie ciągłe.
Dla każdego $i$ przeciwobrazy zbiorów podbazowych przez
$T_i$ są otwarte w $X_i$; ich skończone przecięcia
i dowolne sumy również. Tak więc $T_i:X_i\to(X,\tau)$
są ciągłe. Każda dopuszczalna topologia jest zawarta
w $\tau$, co dowodzi maksymalności. W szczególności
w tej konstrukcji nie trzeba zakładać, że obrazy $T_i$
generują algebraicznie $X$.

Jeśli $S:X\to Z$ jest liniowe, a wszystkie $S\circ T_i$
są ciągłe, to topologia początkowa $\sigma_S$ na $X$
indukowana przez $S$ jest topologią liniowo-topologiczną
(twierdzenie~\ref{thm:initial-tvs}); wszystkie $T_i$
są ciągłe do $(X,\sigma_S)$, a zatem $\sigma_S\subseteq\tau$.
Z definicji $\sigma_S$ wynika, że $S:(X,\tau)\to Z$
jest ciągłe. W drugą stronę wystarczy składać odwzorowania
ciągłe.
\end{proof}

\begin{uwaga}[Wersja lokalnie wypukła]
Jeżeli wszystkie $X_i$ są lokalnie wypukłe, można
analogicznie rozważać \emph{najsilniejszą lokalnie wypukłą}
topologię na $X$, przy której $T_i$ są ciągłe.
Będzie to szczególnie ważne dla przestrzeni funkcji
testowych $\mathcal D(\Omega)$; tam użyjemy granicy
induktywnej w kategorii przestrzeni lokalnie wypukłych.
Nie należy bez dodatkowych argumentów utożsamiać
tej topologii z końcową topologią zwykłych przestrzeni
topologicznych.
\end{uwaga}

\subsection{Sumy proste: produkt skończony i nieskończony}

\begin{definicja}[Algebraiczna suma prosta]
Dla rodziny przestrzeni liniowych $(X_i)_{i\in I}$
przez $\bigoplus_{i\in I}X_i$ oznaczamy przestrzeń
rodzin $(x_i)\in\prod_iX_i$, w których tylko skończenie
wiele współrzędnych jest niezerowych. Naturalne włożenia
$j_i:X_i\to\bigoplus_iX_i$ umieszczają element w $i$-tej
współrzędnej, a na pozostałych wpisują zero.
\end{definicja}

\begin{twierdzenie}[Skończone sumy proste]
Niech $X_1,\dots,X_n$ będą przestrzeniami liniowo-topologicznymi.
Na $X_1\oplus\dots\oplus X_n$ utożsamionym algebraicznie
z $X_1\times\dots\times X_n$ topologia produktowa jest
końcową topologią liniowo-topologiczną względem włożeń $j_i$.
\end{twierdzenie}
\begin{proof}
Włożenia $j_i$ są ciągłe w topologii produktowej.
Jeśli $Z$ jest TVS i $S:X_1\oplus\dots\oplus X_n\to Z$
jest liniowe, a każde $S\circ j_i$ jest ciągłe,
to dla $(x_1,\dots,x_n)$ mamy
$S(x_1,\dots,x_n)=\sum_{i=1}^n(S\circ j_i)(x_i)$.
Skończona suma ciągłych odwzorowań jest ciągła.
Z uniwersalnej własności końcowej topologii
liniowo-topologicznej (twierdzenie~\ref{thm:final-tvs})
wynika teza.
\end{proof}

\begin{przyklad}[Przestrzeń ciągów o skończonym nośniku]\label{ex:direct-sum}
Niech $\varphi=\bigoplus_{n\geq1}\K e_n$ oznacza
przestrzeń ciągów o skończonym nośniku. Topologia
indukowana z produktu $\K^{\N}$ i \emph{topologia
lokalnie wypukłej sumy prostej} są różne.
Dla jasności: \emph{seminorma} to funkcja $p:X\to[0,\infty)$
spełniająca $p(\lambda x)=|\lambda|p(x)$ i
$p(x+y)\leq p(x)+p(y)$; może przyjmować wartość zero
na niezerowych wektorach. Rodzina seminorm definiuje topologię,
w której bazę otoczeń zera stanowią skończone przecięcia
zbiorów $\{x:p(x)<\varepsilon\}$, $\varepsilon>0$.
Jest to topologia \emph{lokalnie wypukła}: zero posiada bazę
wypukłych otoczeń. Teorię seminorm rozwiniemy w \hyperref[sec:lokalnie-wypukle]{rozdziale 5}.

Na $\varphi$ rozważmy rodzinę seminorm
\[
 p_a(x)=\sum_{n\geq1}a_n|x_n|,
 \qquad a=(a_n)_{n\geq1}\in[0,\infty)^{\N}.
\]
Każda suma jest skończona dla ustalonego $x\in\varphi$.
Wszystkie $p_a$ są seminormami, a ich rodzina
rozdziela punkty, więc określają lokalnie wypukłą
topologię Hausdorffa. Niech $Z$ będzie przestrzenią
liniowo-topologiczną z topologią zadaną przez rodzinę
seminorm. Dla liniowego $S:\varphi\to Z$ oraz
dowolnej seminormy $p$ z tej rodziny zachodzi
\[
 p(Sx)\leq\sum_{n\geq1}|x_n|\,p(S e_n)
         =p_a(x),\qquad a_n=p(S e_n).
\]
Wynika stąd ciągłość każdego liniowego $S:\varphi\to Z$:
dla dowolnego $p$ i $\varepsilon>0$ otwarte otoczenie zera
$\{x\in\varphi:p_a(x)<\varepsilon\}$ zawiera się w
$S^{-1}(\{z\in Z:p(z)<\varepsilon\})$.
Ciągłość $S$ w zerze, a następnie w każdym punkcie,
wynika z opisu topologii $Z$ przez seminormy i liniowości $S$.
W \hyperref[sec:lokalnie-wypukle]{rozdziale 5} udowodnimy, że topologie wszystkich
przestrzeni lokalnie wypukłych można opisać seminormami;
wtedy powyższy argument da pełną własność uniwersalną
sumy prostej \emph{w kategorii przestrzeni lokalnie wypukłych}.
Nie dowodzi on analogicznej własności względem wszystkich
przestrzeni liniowo-topologicznych bez lokalnej wypukłości.

W topologii produktowej odziedziczonej z $\K^{\N}$
ciąg $e_n$ zbiega do zera. W topologii sumy prostej
nie zbiega: $p_{(1,1,\dots)}(e_n)=1$ dla każdego $n$.
To kolejny przykład, w którym ten sam zbiór liniowy
niesie dwie naturalne, ale istotnie różne topologie.
\end{przyklad}

\subsection{Iloraz Hausdorffa dowolnej TVS}

\begin{twierdzenie}[Hausdorffyzacja]\label{thm:hausdorffization}
Niech $X$ będzie przestrzenią liniowo-topologiczną
(niekoniecznie Hausdorffa), a $N=\overline{\{0\}}$.
Wówczas $N$ jest domkniętą podprzestrzenią, $X/N$
jest przestrzenią liniowo-topologiczną Hausdorffa
oraz dla każdego ciągłego liniowego $T:X\to Y$,
gdzie $Y$ jest TVS Hausdorffa, istnieje dokładnie
jedno ciągłe liniowe $\widetilde T:X/N\to Y$
spełniające $T=\widetilde T\circ q$.
\end{twierdzenie}
\begin{proof}
Domkniętość i liniowość $N$ wykazano w \hyperref[sec:tvs]{rozdziale 1}. Wobec
\ref{thm:quotient-tvs} iloraz $X/N$ jest Hausdorffa.
Ponieważ $T$ jest ciągłe i $Y$ jest Hausdorffa,
$\ker T=T^{-1}(\{0\})$ jest domknięte i zawiera zero,
a więc $N\subseteq\ker T$. Z własności uniwersalnej
ilorazu (twierdzenie~\ref{thm:quotient-universal})
istnieje jednoznaczne ciągłe $\widetilde T$.
\end{proof}

\subsection{Podsumowanie: kiedy topologia jest właściwa?}

W konstrukcjach tego rozdziału występują cztery pojęcia:
\begin{enumerate}[label=\textup{(\roman*)}]
\item topologia podprzestrzeni jest \emph{najsłabszą} topologią,
      przy której włożenie do większej przestrzeni jest ciągłe;
\item topologia produktowa jest \emph{najsłabszą} topologią,
      przy której wszystkie rzuty współrzędne są ciągłe;
\item topologia ilorazowa jest \emph{najsilniejszą} topologią,
      przy której rzutowanie ilorazowe jest ciągłe; w tym
      przypadku pozostaje topologią przestrzeni liniowo-topologicznej;
\item ogólna topologia końcowa dla rodziny odwzorowań liniowych
      nie musi być zgodna z działaniami liniowymi; trzeba
      rozróżniać wersję topologiczną, liniowo-topologiczną
      i lokalnie wypukłą.
\end{enumerate}

Warto zapamiętać dwie przeszkody: iloraz przez podprzestrzeń
niedomkniętą nie jest Hausdorffa, a ciągła liniowa bijekcja
nie musi być izomorfizmem topologicznym.
W następnym rozdziale zajmiemy się sieciami, strukturą jednostajną
przestrzeni liniowo-topologicznych, zupełnością oraz uzupełnieniami.
